\documentclass[aspectratio=169,10pt,spanish]{beamer}

\input{../../../_preambulo_beamer.tex}
\input{../../../_comandos_beamer.tex}

\selectlanguage{spanish}

\title{MA1001B - Modelación Estadística para la Toma de Decisiones}
\subtitle{Lección 29: Intervalo de confianza para una proporción poblacional}
\author{}
\institute{Tecnol\'ogico de Monterrey}
\date{}

\begin{document}

\begin{frame}[plain]
  \vspace{1.2cm}
  {\Large\bfseries MA1001B - Modelación Estadística para la Toma de Decisiones\par}
  \vspace{0.7cm}
  {\large Lección 29: Intervalo de confianza para una proporción poblacional\par}
  \vspace{0.7cm}
  {\color{TecRojo}\rule{\textwidth}{1.2pt}}
  \vspace{0.8cm}

  Facultad MA1001B\\[0.3cm]
  Periodo\\[0.3cm]
  Tecnol\'ogico de Monterrey
\end{frame}

\begin{frame}{La pregunta de la tasa binaria}
  \begin{alertblock}{Pregunta guía}
    ¿Cómo se construye un intervalo al 95\% para una proporción desconocida
    a partir de $x$ éxitos en $n$ ensayos independientes?
  \end{alertblock}

  \vspace{0.6em}
  Al final de esta lección, usted puede:
  \begin{itemize}
    \item calcular $\hat{p}=x/n$;
    \item verificar $n\hat{p}$ y $n(1-\hat{p})$ antes de usar $z$;
    \item formar el intervalo de Wald $\hat{p}\pm 1.96\sqrt{\hat{p}(1-\hat{p})/n}$;
    \item explicar por qué un evento raro con $x$ pequeña hace engañoso a Wald.
  \end{itemize}

  \vfill
  \scriptsize Todos los conteos de tickets usados hoy son sintéticos. No se usan datos reales de empresa.
\end{frame}

\begin{frame}{Una instantánea sintética de primer contacto}
  \small
  \begin{center}
    \begin{tabular}{lr}
      \toprule
      \textbf{Cantidad} & \textbf{Valor} \\
      \midrule
      Tickets $n$ & 150 \\
      Resoluciones en primer contacto $x$ & 27 \\
      No primer contacto & 123 \\
      Proporción muestral $\hat{p}$ & $27/150=0.180000$ \\
      \bottomrule
    \end{tabular}
  \end{center}

  \vspace{0.5em}
  \begin{exampleblock}{Verificación éxito-fracaso}
    $n\hat{p}=27\ge 5$ y $n(1-\hat{p})=123\ge 5$. La aproximación normal de
    $\hat{p}$ se trata como usable para esta instantánea.
  \end{exampleblock}
\end{frame}

\begin{frame}[fragile]{Paso 1: Proporción muestral}
  \begin{columns}[T]
    \begin{column}{0.42\textwidth}
      \small
      \begin{lstlisting}
phat = x / n
n_phat = n * phat
n_qhat = n * (1 - phat)
      \end{lstlisting}

      \begin{block}{Salida verificada}
        $\hat{p}=0.180000$\\
        $n\hat{p}=27$\\
        $n(1-\hat{p})=123$
      \end{block}
    \end{column}
    \begin{column}{0.58\textwidth}
      \centering
      \includegraphics[width=\textwidth]{../figuras/prop_ci_01_sample_proportion.png}
      \vspace{0.2em}
      \scriptsize Conteos de primer contacto del Paso 1
    \end{column}
  \end{columns}
\end{frame}

\begin{frame}{Intervalo de Wald después de que la verificación pasa}
  \small
  Error estándar y margen al 95\%:
  \[
    \mathrm{EE}=\sqrt{\frac{0.18\times 0.82}{150}}=0.031369,
  \]
  \[
    E=1.96\times 0.031369=0.061483.
  \]
  Intervalo:
  \[
    \hat{p}\pm E=[0.118517, 0.241483].
  \]
  El intervalo estima la tasa poblacional desconocida $p$, no el siguiente
  ticket.
\end{frame}

\begin{frame}[fragile]{Paso 2: El intervalo de Wald al 95\%}
  \begin{columns}[T]
    \begin{column}{0.40\textwidth}
      \small
      \begin{lstlisting}
se = np.sqrt(
    phat * (1 - phat) / n
)
lower = phat - 1.96 * se
      \end{lstlisting}

      \begin{block}{Salida verificada}
        $\mathrm{EE}=0.031369$\\
        Margen $=0.061483$\\
        IC $[0.118517, 0.241483]$
      \end{block}
    \end{column}
    \begin{column}{0.60\textwidth}
      \centering
      \includegraphics[width=\textwidth]{../figuras/prop_ci_02_wald_interval.png}
      \vspace{0.2em}
      \scriptsize Intervalo de Wald del Paso 2
    \end{column}
  \end{columns}
\end{frame}

\begin{frame}{Paso 3: Los eventos raros hacen engañoso a Wald}
  \begin{columns}[T]
    \begin{column}{0.42\textwidth}
      \small
      Ahora $x=3$, $n=150$, $\hat{p}=0.020000$.

      \vspace{0.4em}
      $n\hat{p}=3<5$. Las condiciones fallan.

      \vspace{0.4em}
      Intervalo de Wald:
      \[
        [-0.002405, 0.042405].
      \]
      El límite inferior es negativo.

      \vspace{0.5em}
      \textbf{Límite:} una proporción no puede ser negativa. Una $x$ pequeña
      hace engañoso el intervalo de Wald.
    \end{column}
    \begin{column}{0.58\textwidth}
      \centering
      \includegraphics[width=\textwidth]{../figuras/prop_ci_03_rare_event_wald.png}
      \vspace{0.2em}
      \scriptsize Diagnóstico Wald versus Wilson del Paso 3
    \end{column}
  \end{columns}
\end{frame}

\begin{frame}{Verificación de conceptos}
  \begin{enumerate}
    \item Calcule $\hat{p}$ a partir de $x=27$ y $n=150$.
    \item ¿Pasan $n\hat{p}=27$ y $n(1-\hat{p})=123$ el umbral 5?
    \item ¿Por qué un límite de Wald negativo es un problema para una
      proporción?
    \item ¿Qué cantidad falla la verificación cuando $x=3$ y $n=150$?
  \end{enumerate}

  \vspace{0.6em}
  \begin{block}{Conexión con la Lección 30}
    La sesión puede continuar directamente con la elección de $n$ de modo
    que un margen de error objetivo sea alcanzable bajo un presupuesto.
  \end{block}

  \vfill
  \scriptsize Fuente: OpenStax, \textit{Introductory Business Statistics 2e},
  Capítulo 8, Sección 8.3. CC BY-NC-SA.
\end{frame}

\appendix



\end{document}
